Formulation of a penetration conditions \hfill[2]

From the conditions specified in the exercise text the particle should be pass in the belt of
magnetic field around quarter of the circle. If this particle leaves the belt the particle path
should be a little longer. It means that the radius of this circle should be larger then the belt
thickness.
% sic: "should be pass", "larger then" as in the original.

Determination of the forces (centrifugal and magnetic) acting upon the particle, writing of the
appropriate formulae \hfill[4]

Comparison of the forces and calculations \hfill[3]

Comparing centrifugal and magnetic force one obtains
\[
e \cdot v \cdot B = M \cdot v^2 / r
\]
So $e \cdot r \cdot B = p$, where $p$ denotes the momentum. From the assumed approximation
$e \cdot r \cdot B = E / c$.

Thus definitively particle energy is described by $E = c \cdot e \cdot r \cdot B$.

Substituting numerical values one obtains $E \approx 3\cdot10^{-8}$ J it means
$2\cdot10^{2}$ GeV.

Answer and conclusion: \hfill[1]
